\documentclass[11pt]{article}
\usepackage[margin=1in]{geometry}
\usepackage{booktabs}
\usepackage{amsmath}
\usepackage{amssymb}
\usepackage{graphicx}
\usepackage{xcolor}
\usepackage{hyperref}
\hypersetup{colorlinks=true, linkcolor=blue, urlcolor=blue, citecolor=blue}
\usepackage{enumitem}
\usepackage{array}

\newcommand{\meas}{\textcolor{blue!60!black}{\small[measured]}}
\newcommand{\modeled}{\textcolor{orange!70!black}{\small[modeled]}}
\newcommand{\cited}{\textcolor{purple!70!black}{\small[cited]}}

\title{\textbf{The Unified Sovereign-Tenant Cloud:\\
Linear Shared-Nothing Scaling via Per-Tenant SQLite\\
and In-Process Composition}}
\author{Hanzo AI Research \\ \textit{Hanzo Industries Inc} \\ \texttt{research@hanzo.ai}}
\date{Measurement pass: 2026-07-08 (host \texttt{evo}, AMD Ryzen AI Max+ 395, Go 1.26.4)}

\begin{document}
\maketitle

\begin{abstract}
Hanzo Cloud runs the entire control- and data-plane of a multi-tenant SaaS as \emph{one}
ZAP-native Go binary (56 in-process subsystems, one runtime) over a \emph{shared-nothing}
storage substrate: every tenant --- org, app, project, or user --- is its own embedded SQLite
file, replicated to durable object storage as per-tenant post-quantum-encrypted LTX segments,
served by a single pinned writer per file, with only the \emph{active} tenants held open in a
bounded per-pod LRU. This paper asks one question and answers it with measurement, not
marketing: \emph{what does a resident tenant actually cost, and how does the design scale?}
We measure the marginal resident-set size (RSS) of opening $N \in \{1,10,100,1000\}$ real
base-schema tenant databases through the exact production open path, on the production driver.
The answer: \textbf{an open, actively-served tenant costs $\approx$\,0.25\,MB of RAM}
(251\,KB measured incremental slope; cgo/SQLite, warm workload) --- roughly $8\times$ below the
2\,MB figure the design was previously budgeted against --- and an \emph{idle} tenant costs
essentially zero live RAM (it is a file on disk plus an encrypted object). The tuned
\texttt{cache\_size} cap is a \emph{ceiling}, not a footprint: for realistic tenants the working
set, not the cap, sets the cost. From this one measured constant we build a capacity model: at
1\% peak concurrency, \textbf{one billion registered tenants} is served by
$\approx$\,103 32\,GB nodes ($\approx$\,2.5\,TB cluster RAM); at 0.1\% concurrency, 11 nodes.
Memory scales with \emph{concurrency}, not with registrations. We contrast this against the
default industry shape --- shared Postgres behind a microservice mesh --- and quantify the
penalties it pays that the unified design does not: a global write bottleneck, a hard
connection ceiling, and an $N$-runtime plus sidecar memory tax we measure at
\textbf{$11$--$32\times$} the unified binary's footprint. We are explicit about the limits:
writer-pinning failover, hot-tenant skew, cross-tenant queries, and the routing directory are
real costs, discussed honestly in \S\ref{sec:limits}. Every number labelled \meas{} is measured
on the named host; \modeled{} and \cited{} numbers are labelled as such; none are invented.
\end{abstract}

\section{Introduction}
\label{sec:intro}
The dominant shape of a multi-tenant SaaS backend is a fleet of stateless microservices in
front of a shared relational database (typically Postgres), with tenancy expressed as a
\texttt{tenant\_id} column and enforced in application code or row-level security. That shape
has two structural costs that grow with the business. First, \emph{writes serialize}: every
tenant's mutations funnel through one primary's write-ahead log, so the write path is a single
global resource no matter how many read replicas are added. Second, \emph{runtimes multiply}:
each microservice is a full process --- its own language runtime, heap, connection pool, and
(in a service mesh) a sidecar proxy --- so the fixed memory cost of the platform is paid $N$
times over, before a single byte of tenant data is stored.

Hanzo Cloud takes the opposite position on both axes, and this paper is the measured defense of
that position. On composition: the 56 subsystems (IAM, KMS, Base, Commerce, AI, o11y, VFS, MQ,
tasks, audit, \dots) are linked into \emph{one} binary and communicate in-process over a
ZAP-native transport, with the same symbols exposed as ZAP-RPC only where process isolation is
a compliance requirement (payments/vault, PCI CDE)~\cite{hip0106,hip0116,hip0125}. The unified
binary's construction is developed in the companion architecture paper~\cite{unifiedbin}; this
paper is its measured memory and scaling counterpart. On storage: there is
no shared database. Each tenant is a sovereign, independently encrypted SQLite
file~\cite{hip0302}; the write path for tenant $t$ is a single writer pinned to $t$, so writes
are \emph{embarrassingly sharded} and there is no global lock, no distributed transaction, and
no per-write consensus round. The claim is that this is simultaneously the more \emph{efficient}
design (less RAM per tenant, far less fixed overhead) and the more \emph{scalable} one (node
count grows linearly in active tenants, with no shared bottleneck to saturate).

Claims about efficiency and scale are cheap; the contribution here is to make them measurable.
We (i) establish the real per-active-tenant memory constant by direct measurement through the
production open path (\S\ref{sec:bench}); (ii) turn that constant into a capacity model from
100 to $10^9$ registered tenants (\S\ref{sec:capacity}); (iii) quantify the memory gap against
the shared-Postgres-plus-mesh alternative (\S\ref{sec:compare}); and (iv) state the scaling
claims as propositions grounded in the measurement, with an honest accounting of where the
design pays (\S\ref{sec:claims}--\S\ref{sec:limits}).

\section{Architecture}
\label{sec:arch}

\paragraph{One binary, 56 in-process subsystems.}
Per HIP-0106~\cite{hip0106}, every Hanzo Go service exposes a top-level \texttt{Mount(app,
deps)} symbol and is compiled into a single \texttt{cloud} binary rather than deployed as an
independent service. Subsystems consume each other through Go interfaces that resolve to
\emph{in-process} implementations by default (the in-process plugin/composition model of
HIP-0116~\cite{hip0116}), falling back to ZAP-RPC only for the small set of subsystems that must
stay in a separate process for compliance isolation. HIP-0125~\cite{hip0125}
eliminates gRPC from the transport entirely in favor of ZAP-native framing, so there is no
per-subsystem gRPC server, no protobuf reflection service, and no sidecar. The consequence for
this paper is that the platform's \emph{fixed} memory cost --- one Go runtime, one scheduler,
one set of shared caches and connection machinery --- is paid \emph{once} per node, not once per
subsystem.

\paragraph{The sovereign tenant: one encrypted SQLite file each.}
Tenancy is a composable hierarchy under one org: an org-wide database, a per-app database, a
per-app-per-project database, and a per-user database~\cite{hip0302}. Each lands on disk at
\texttt{\{DataDir\}/orgs/\{orgSlug\}/\{service\}.db} and is opened through exactly one code path
--- the canonical Hanzo SQLite driver with a fixed pragma battery (\S\ref{sec:bench}). There is
no shared schema and no cross-tenant table: tenant isolation is a \emph{filesystem} boundary,
not a \texttt{WHERE tenant\_id = ?} predicate, which removes an entire class of
tenant-bleed bugs by construction~\cite{hip0118}.

\paragraph{The pinned single writer.}
A single tenant is served by exactly one pod at a time via gateway-side sticky affinity
(consistent hash on \texttt{X-Org-Id}). Because SQLite is a single-writer store and each tenant
maps to one file served by one pod, the writer for tenant $t$ is globally unique without any
coordination protocol. Reads may be served by read-only replica pods that hydrate from the
replication stream and fail \emph{closed} on any mutating verb, so correctness never rides on
routing being perfect. This is the mechanism behind the linear-scaling claim: distinct tenants
share no write path, so adding tenants adds independent writers rather than contending for one.

\paragraph{Durability: encrypted replication, not a shared database.}
Per HIP-0302~\cite{hip0302}, each tenant's SQLite file is continuously replicated as LTX
segments through \texttt{encreplica}, which post-quantum-\texttt{age}-encrypts every segment
with a \emph{per-tenant} key before it touches durable storage (the HIP-0107 replication
substrate~\cite{hip0107}; local FS for dev, S3 or Hanzo VFS in production~\cite{hip0032}). The
storage and encryption architecture is developed in the companion post-quantum object-store
paper~\cite{pqstore}; this paper concerns its memory and scaling behavior. At-rest encryption
therefore lives on the replication path,
not in the resident page cache: the open handle is plaintext SQLite (fast), and the durable
copy is ciphertext keyed to the tenant. This distinction matters for the memory model --- the
crypto cost is transient (a buffer during flush), not a standing per-connection tax.

\paragraph{Active-set residency.}
Only \emph{active} tenants are held open. The per-pod store keeps open SQLite handles in a
bounded LRU (default cap 1000 handles/pod, tunable) with a 5-minute idle reaper; on eviction it
checkpoints the WAL, closes the handle, and flushes the encrypted bytes to the object store. A
cold tenant occupies no RAM --- it is an encrypted object plus, at most, a file in a local
cache. This is the architectural reason the resident memory of a node tracks \emph{concurrency}
rather than \emph{registrations}, which \S\ref{sec:bench} then quantifies.

\section{Benchmark: the real per-tenant memory constant}
\label{sec:bench}

\paragraph{Question.}
What is the marginal resident memory of one open, actively-served tenant database, measured
through the exact production open path, on the driver the production binary actually links?

\paragraph{Method.}
We open tenant databases with the identical call the production store uses,
\begin{center}
\texttt{dbx.Open("sqlite", sqlite.PragmaDSN(path, sqlite.DefaultPragmas))},
\end{center}
where \texttt{DefaultPragmas} is the canonical Hanzo battery: \texttt{busy\_timeout=10000},
\texttt{journal\_mode=WAL}, \texttt{synchronous=NORMAL}, \texttt{foreign\_keys=ON},
\texttt{temp\_store=MEMORY}, \texttt{cache\_size=-32000} (a 32\,MiB page-cache \emph{cap} per
connection). The seed is the checked-in base test database
(\texttt{base/tests/data/data.db}, 344{,}064\,B data + 49{,}152\,B auxiliary $\approx$ 0.38\,MB),
a realistic small tenant carrying the base system collections plus a little data. For
$N \in \{1,10,100,1000\}$ we copy the seed to a fresh path, open it, run a workload, and keep
the handle resident; at each checkpoint we force a Go GC and \texttt{FreeOSMemory} and read
\texttt{VmRSS} from \texttt{/proc/self/status}. We report the \emph{incremental slope}
$\Delta\text{RSS}/\Delta N$ between checkpoints, which cancels the fixed process baseline and
yields the true per-tenant constant. RSS (not Go heap statistics) is the honest measure because
the cgo SQLite page cache is allocated in the C heap, invisible to Go's allocator. Two workloads:
\emph{cold} (open + a single \texttt{SELECT 1}) and \emph{warm} (touch every table's pages via
\texttt{count(*)} then perform a small write, materializing the WAL, the \texttt{-shm} index,
and the pinned writer connection). The production driver is cgo/\texttt{mattn} (the \texttt{cloud}
binary links it as the active \texttt{"sqlite"} driver); we also measure the pure-Go
\texttt{modernc} backend used on \texttt{!cgo} builds. Host: \texttt{evo}, AMD Ryzen AI Max+ 395
(32 cores), 62.4\,GB, Linux 6.17.0-35, Go 1.26.4.

\paragraph{Result: $\approx$0.25\,MB per active tenant.}
Table~\ref{tab:perdb} gives the measured constants. On the production driver, a warm resident
tenant costs a measured \textbf{251\,KB} ($\approx$0.245\,MB) at the asymptotic
(100\,$\rightarrow$\,1000) slope; a cold-but-open tenant costs 235\,KB. The pure-Go backend is
$\approx$1.5$\times$ heavier at 392\,KB. The single most important finding is in the
\texttt{cache\_size} column: sweeping the cap across $\{-256, -1024, -2000, -32000\}$
(256\,KiB to 32\,MiB) moves the small-tenant constant by \emph{less than 0.2\%}
(251.3--251.7\,KB). The 32\,MiB cap is never approached, because the tenant's working set
($\approx$0.25\,MB) --- not the cap --- sets the footprint. The four cache settings are in effect
four independent repeats at equal working set, and their agreement to within 0.4\,KB is our
repeatability evidence.

\begin{table}[h]
\centering
\caption{Measured marginal RSS per open, resident tenant (base-schema seed, 0.38\,MB on disk).
Incremental slope is $\Delta\text{RSS}/\Delta N$ over $100\rightarrow1000$ handles; baseline
(runtime + driver, zero tenants) is 7.0--7.2\,MB (cgo), 6.3--6.5\,MB (modernc). \meas{}}
\label{tab:perdb}
\begin{tabular}{llrr}
\toprule
Driver & Workload & \texttt{cache\_size} & Per-tenant RSS \\
\midrule
mattn / SQLite (cgo, \emph{production}) & warm & $-32000$ (32\,MiB) & \textbf{251.3\,KB} (0.245\,MB) \\
mattn / SQLite (cgo, \emph{production}) & warm & $-2000$ (2\,MiB)   & 251.4\,KB \\
mattn / SQLite (cgo, \emph{production}) & warm & $-1024$ (1\,MiB)   & 251.3\,KB \\
mattn / SQLite (cgo, \emph{production}) & warm & $-256$ (256\,KiB)  & 251.4\,KB \\
mattn / SQLite (cgo, \emph{production}) & cold & $-32000$           & 235.5\,KB (0.230\,MB) \\
modernc (pure-Go, \texttt{!cgo})        & warm & $-32000$           & 392.0\,KB (0.383\,MB) \\
\bottomrule
\end{tabular}
\end{table}

\paragraph{The density knob is real --- for \emph{hot, large} tenants only.}
That \texttt{cache\_size} is inert for small tenants is not the whole story. We re-ran the sweep
against a synthetic 25\,MB tenant under a full-scan warm workload (Table~\ref{tab:knob}). Here the
cap binds: at $-32000$ a fully-scanned hot tenant fills its page cache to the working set and
costs 25.5\,MB; at $-1024$ it is bounded to 1.18\,MB; at $-256$ to 0.41\,MB. So the cap is a
per-tenant memory \emph{ceiling} operators can dial to trade page-cache hit rate for density.
The realistic operating point is Table~\ref{tab:perdb} (working-set-bound, $\approx$0.25\,MB);
Table~\ref{tab:knob} is the worst case (analytic tenant scanning its whole database every
request) and shows it remains bounded and tunable rather than unbounded.

\begin{table}[h]
\centering
\caption{Density knob: marginal RSS for a \emph{hot} 25\,MB tenant under a full-table-scan
workload, as a function of the \texttt{cache\_size} cap. \meas{}}
\label{tab:knob}
\begin{tabular}{lrr}
\toprule
\texttt{cache\_size} cap & Per-tenant RSS & Effect \\
\midrule
$-32000$ (32\,MiB) & 25.5\,MB & cache fills to working set (cap not reached) \\
$-1024$ (1\,MiB)   & 1.18\,MB & bounded to $\approx$1\,MiB regardless of DB size \\
$-256$ (256\,KiB)  & 0.41\,MB & bounded to $\approx$256\,KiB + fixed overhead \\
\bottomrule
\end{tabular}
\end{table}

\paragraph{Idle tenants cost $\approx$0.}
We opened and exercised 1000 tenants, then closed every handle. On the pure-Go backend --- whose
memory frees to the Go heap, which \texttt{FreeOSMemory} returns to the OS --- RSS fell from
392\,MB back to \textbf{11.2\,MB}, i.e.\ to within 4.9\,MB of baseline over 1000 closed tenants
($\approx$5\,KB/tenant of retained runtime bookkeeping). This is the direct measurement behind
``an idle tenant is free'': closing a handle releases essentially all of its live memory. On the
cgo backend RSS settled at 221\,MB after close; that residual is \emph{not} live memory but
glibc's retention of freed C arenas (the allocator does not \texttt{munmap} them back to the
kernel). We confirmed it is reused, not leaked: closing 1000 tenants and opening 1000 fresh ones
grew RSS by only $\approx$110\,MB on top of the 222\,MB retained (vs.\ 247\,MB for a cold start),
so the freed arenas absorb subsequent opens. The operational takeaway: a pod's steady-state RSS
tracks its \emph{high-water active set}, and idle registrations contribute no live pages.

\paragraph{On-disk and base-node constants.}
A realistic tenant occupies $\approx$0.38\,MB on disk (measured seed), which at rest lives as
per-tenant-encrypted LTX in object storage~\cite{hip0032,hip0302}, not on node-local disk. The
fixed per-node cost --- the \texttt{cloud} binary with all 56 subsystems and one Go runtime --- is
a working-set RSS of \textbf{154--225\,MB} from prior production measurement~\cite{edgepaper};
we use 190\,MB as the point estimate in the model below and note it is \cited{} (the cluster was
not reachable from the measurement host on this pass to re-sample it live). For comparison, a
\emph{bare} Go HTTP service (net/http, one handler, no framework) measured \meas{} 7.2\,MB on the
same host --- the floor a per-service split cannot go below, used in \S\ref{sec:compare}.

\section{The capacity model}
\label{sec:capacity}
The measured constant makes the fleet size a closed-form function of registrations and
concurrency. Let $R$ be registered tenants, $c$ the peak concurrency fraction, so active
tenants $A = cR$. Let $m$ be the measured per-active-tenant RAM (MB), $b$ the per-node base RSS
(GB), and $U$ the usable RAM per node (GB) after reserving for the OS, kubelet, and burst
headroom. Then
\begin{equation}
\text{nodes} \;=\; \max\!\Big(3,\; \Big\lceil \frac{A \cdot m}{(U - b)\cdot 1024} \Big\rceil\Big),
\qquad
\text{cluster RAM} \;=\; \text{nodes}\cdot b \;+\; \frac{A\cdot m}{1024}\ \text{GB}.
\label{eq:model}
\end{equation}
We use 32\,GB nodes with $U = 24$\,GB (75\% usable), $b = 0.19$\,GB, and the measured
$m = 0.25$\,MB, giving a per-node density of $\approx$97{,}500 active tenants; we also report the
conservative $m = 1.0$\,MB envelope (larger/hotter tenants, multiple resident connections,
extra headroom), which yields $\approx$24{,}400 active/node. Disk is $R \times 0.38$\,MB total,
object-tiered. Table~\ref{tab:cap} gives the model at 1\% peak concurrency;
Table~\ref{tab:conc} shows concurrency sensitivity at the extremes.

\begin{table}[h]
\centering
\caption{Capacity at \textbf{1\% peak concurrency}, 32\,GB nodes, from Eq.~\ref{eq:model} with the
\meas{} $m=0.25$\,MB constant (and the \modeled{} $m=1.0$\,MB conservative envelope). ``Cluster
RAM'' is consumed RAM per Eq.~\ref{eq:model}, not provisioned. Node counts are floored at 3 for HA.}
\label{tab:cap}
\resizebox{\textwidth}{!}{%
\begin{tabular}{r r r r r r}
\toprule
Registered $R$ & Active @1\% & Nodes ($m{=}0.25$) & Cluster RAM & Nodes ($m{=}1.0$) & Disk (total, tiered) \\
\midrule
100            & 1          & 3   & 0.57\,GB  & 3    & 40\,MB  \\
1{,}000        & 10         & 3   & 0.57\,GB  & 3    & 400\,MB \\
10{,}000       & 100        & 3   & 0.59\,GB  & 3    & 4\,GB   \\
100{,}000      & 1{,}000    & 3   & 0.81\,GB  & 3    & 40\,GB  \\
1{,}000{,}000  & 10{,}000   & 3   & 3.01\,GB  & 3    & 400\,GB \\
10{,}000{,}000 & 100{,}000  & 3   & 24.98\,GB & 5    & 4\,TB   \\
100{,}000{,}000& 1{,}000{,}000 & 11 & 246.2\,GB & 42 & 40\,TB  \\
1{,}000{,}000{,}000 & 10{,}000{,}000 & \textbf{103} & \textbf{2.46\,TB} & 411 & 400\,TB \\
\bottomrule
\end{tabular}}
\end{table}

\begin{table}[h]
\centering
\caption{Concurrency sensitivity for $R = 10^9$ registered tenants ($m=0.25$\,MB, 32\,GB nodes).
Memory scales with \emph{concurrency}, not registrations: a 50$\times$ swing in peak concurrency
is a 50$\times$ swing in fleet size, while $R$ is held fixed. \meas{}-grounded.}
\label{tab:conc}
\begin{tabular}{r r r r}
\toprule
Peak concurrency & Active tenants & Nodes & Cluster RAM \\
\midrule
0.1\% & 1{,}000{,}000  & 11  & 246\,GB \\
1.0\% & 10{,}000{,}000 & 103 & 2.46\,TB \\
5.0\% & 50{,}000{,}000 & 513 & 12.3\,TB \\
\bottomrule
\end{tabular}
\end{table}

The shape of Table~\ref{tab:cap} is the thesis in numbers. Up to 10\,million registered tenants
at 1\% concurrency, the fleet is the 3-node HA floor --- the platform is base-RSS-bound, not
tenant-bound, and the tenants are ``free'' relative to the fixed cost of running at all. Only
past $\sim$$10^8$ registrations does tenant memory dominate, and there it grows \emph{linearly}:
$10^9$ tenants is $103$ nodes, $10^{10}$ would be $\sim$1030. There is no knee, no bottleneck
tier, and no rewrite between the 3-node and the 103-node regime --- the same binary and the same
per-tenant file, more of them.

\section{Comparative analysis: shared Postgres + microservice mesh}
\label{sec:compare}
We compare against the default industry shape on the three axes where the unified design's
structure diverges. The comparison is deliberately conservative: where we estimate the
alternative we anchor to a measured floor and label the estimate \modeled{}.

\paragraph{(1) The write path: sharded vs.\ global.}
In shared Postgres, every tenant's writes commit through one primary's WAL. Read replicas scale
reads but not writes; the write ceiling is one node's fsync throughput, and it is \emph{shared}
across all tenants, so a write-heavy tenant degrades every other tenant's write latency. Scaling
writes past one primary requires sharding (Citus, Vitess-style) or multi-primary consensus
(Raft/Paxos per write) --- either a re-architecture or a per-commit coordination round. In the
unified design the write path for tenant $t$ is a private SQLite writer; $N$ tenants are $N$
independent writers with \emph{no shared lock and no consensus round}. Aggregate write
throughput is the sum of per-tenant throughputs and grows linearly with nodes. The cost paid for
this is that a single tenant's writes are bounded by one file's single-writer throughput
(\S\ref{sec:limits}); the benefit is that there is no global write resource to saturate.

\paragraph{(2) The connection ceiling.}
A Postgres primary allocates a backend process per connection ($\sim$5--10\,MB each), so
practical \texttt{max\_connections} is in the hundreds; thousands of tenant-affine connections
require a pooler (PgBouncer) and still contend for the one primary. Per-tenant SQLite has no
connection ceiling: a ``connection'' is an in-process handle into a memory-mapped file, and the
active-set LRU (\S\ref{sec:arch}) bounds the resident count per pod directly. There is no network
round trip and no server-side backend process on the data path at all.

\paragraph{(3) The $N$-runtime + sidecar memory tax.}
This is the axis we can quantify most directly. The unified binary runs all 56 subsystems in one
Go runtime at a measured/cited base of 190\,MB; at 3 replicas for HA that is 570\,MB of
control-plane RAM before tenant data. A microservice mesh runs each subsystem as its own process
with its own runtime heap and, under a service mesh, a sidecar proxy. Even at the \emph{measured
floor} of a bare Go service (7.2\,MB \meas{}) plus a single Istio/Envoy sidecar
($\approx$50\,MB \cited~\cite{istio}) at 2 replicas, 56 subsystems cost 6.4\,GB; at a realistic
60\,MB/service (\modeled{}: framework + OTEL SDK + DB pool + business logic) plus 50\,MB sidecar,
12.3\,GB. Table~\ref{tab:mesh} gives the multiples: the mesh spends
\textbf{$11$--$32\times$} the unified binary's fixed memory, and adds one network hop (plus TLS
and serialization) to every inter-subsystem call that the unified design makes as an in-process
function call.

\begin{table}[h]
\centering
\caption{Fixed control-plane memory: unified binary vs.\ 56-subsystem microservice mesh, before
any tenant data. Unified base 190\,MB is \meas{}/\cited{}; bare-Go 7.2\,MB is \meas{}; per-service
framework RSS and sidecar RSS are \modeled{}/\cited{}. Multiple is mesh\,$\div$\,unified.}
\label{tab:mesh}
\resizebox{\textwidth}{!}{%
\begin{tabular}{l r r r}
\toprule
Deployment & Per-service & Fixed RAM & $\times$ unified (3-rep, 570\,MB) \\
\midrule
Unified binary, 3 replicas          & --                & 570\,MB   & 1.0$\times$ \\
Mesh, floor (7.2\,MB svc + 50\,MB sidecar, 2 rep) & 57.2\,MB & 6.4\,GB & 11.2$\times$ \\
Mesh, realistic (60\,MB svc + 50\,MB sidecar, 2 rep) & 110\,MB & 12.3\,GB & 21.6$\times$ \\
Mesh, heavy (120\,MB svc + 70\,MB sidecar, 2 rep) & 190\,MB & 21.3\,GB & 37.3$\times$ \\
\bottomrule
\end{tabular}}
\end{table}

Composed with the capacity model, the gap compounds: the mesh's 6--21\,GB fixed tax is paid on
\emph{every node}, so at the 3-node HA floor --- where the unified platform serves up to
$10^7$ tenants in $<$1\,GB of control-plane RAM (Table~\ref{tab:cap}) --- the mesh alternative
already needs 19--64\,GB just to exist.

\section{Claims, grounded in the measurement}
\label{sec:claims}
We state the design's guarantees as propositions, each tied to a measured or structural fact
rather than an assertion.

\paragraph{Proposition 1 (memory $\propto$ concurrency, not registrations).}
Resident memory is $\text{nodes}\cdot b + A\cdot m$ with $A = cR$ (Eq.~\ref{eq:model}). Because
an idle tenant contributes $\approx$0 live RAM (measured: 1000 closed handles return to within
4.9\,MB of baseline, \S\ref{sec:bench}) and only the active-set LRU is resident, RAM is a
function of $A$, and $R$ enters only through $A = cR$. Doubling registrations at fixed
concurrency does not change resident memory; it changes only object-storage bytes.
Table~\ref{tab:conc} is the empirical shape.

\paragraph{Proposition 2 (linear node scaling from shared-nothing writers).}
Distinct tenants share no write path: each is a private single-writer SQLite file pinned to one
pod (\S\ref{sec:arch}). There is no global lock, no distributed transaction, and no per-write
consensus, so adding a node adds independent write and read capacity with no coordination
overhead that grows with the fleet. Fleet size is therefore $\Theta(A)$ ---
$103$ nodes at $10^7$ active, $\sim$1030 at $10^8$ --- with a constant factor set by the measured
$m$. This is linear, not merely sublinear-with-a-bottleneck.

\paragraph{Proposition 3 (the one-binary memory advantage is multiplicative).}
Composing $S$ subsystems into one runtime pays the fixed runtime + sidecar cost once rather than
$S$ times. Measured against a 56-subsystem mesh, the fixed-memory ratio is $11$--$32\times$
(Table~\ref{tab:mesh}), and it is paid per node, so it multiplies the per-node base $b$ in
Eq.~\ref{eq:model}. Lowering $b$ is what keeps small deployments in the sub-gigabyte regime in
Table~\ref{tab:cap}.

\section{Limitations}
\label{sec:limits}
The rigor is the persuasion, so the costs are stated as plainly as the benefits.

\begin{itemize}[leftmargin=1.4em]
\item \textbf{Writer-pinning failover.} One writer per tenant means a pod loss requires
re-pinning that tenant to a new pod and replaying the tail of the encrypted replication stream
before writes resume. HIP-0302's consistency model is at-most-once under rebalance: a short
single-writer overlap window can exist during an HPA move, and operators must drain before
scaling. There is a bounded write-unavailability window per affected tenant on failover; the
CAS (ETag If-Match) follow-up is required for generation-checked writes across a rebalance.

\item \textbf{Hot-tenant skew.} Capacity is modeled on an \emph{average} tenant. One tenant that
is simultaneously large and hot can, at the default 32\,MiB cap, cost up to 25.5\,MB resident
(Table~\ref{tab:knob}) and is bounded to one pod's single-writer throughput. The mitigations ---
a smaller per-tenant \texttt{cache\_size}, splitting a heavy org across the app/project/user
hierarchy, and per-tenant rate ceilings --- are real levers but they are operational work, not
free. The model's $m$ should be chosen for the p99 tenant, not the mean, in skewed populations.

\item \textbf{Cross-tenant queries.} A filesystem isolation boundary is exactly the wrong shape
for a query that must span tenants (fleet-wide analytics, admin search, billing rollups). These
require fan-out plus aggregation, or a separate columnar sink fed from the replication stream ---
neither is a single SQL statement. The design optimizes the 99\% per-tenant path at the cost of
the cross-tenant path.

\item \textbf{The routing directory.} Sticky single-writer routing needs a consistent, durable
map from tenant to pod. That directory is itself a piece of shared state that must be highly
available and consistent; it is small and changes slowly, but it is a coordination point the
pure shared-nothing story would prefer not to have.

\item \textbf{Measurement scope.} The per-tenant constant is measured on one host, one seed
schema, single-stream warm/cold workloads, with a bounded 2-connection pool. It does not sweep
concurrent query load per tenant (which materializes more connections, each with its own page
cache), nor multi-hour steady-state fragmentation. The 190\,MB base RSS is \cited{} from prior
production sampling, not re-measured on this pass. These are the honest edges of the numbers.
\end{itemize}

\section{Conclusion}
\label{sec:conclusion}
The unified sovereign-tenant design is efficient and scalable for the same structural reason:
it refuses to share the two resources that normally become bottlenecks. It does not share a
write path --- each tenant is a private single-writer SQLite file --- so writes shard perfectly
and node count grows linearly in active tenants. It does not share $N$ runtimes --- 56 subsystems
live in one binary --- so the fixed platform cost is paid once, at $11$--$32\times$ less memory
than the equivalent mesh. The measured per-active-tenant cost is $\approx$0.25\,MB, an idle
tenant is free, and the \texttt{cache\_size} cap is a tunable ceiling rather than a footprint.
The capacity model that follows from that one constant puts a billion registered tenants on
$\sim$100 commodity nodes at 1\% concurrency, and on 11 at 0.1\% --- with no bottleneck tier and
no rewrite between the small and the large regime. The design is not free of cost: writer
failover, hot-tenant skew, and cross-tenant queries are genuine prices, stated in
\S\ref{sec:limits}. But on the axis it is built for --- serving an enormous number of
independent tenants cheaply and scaling their concurrency linearly --- the measurement supports
the architecture.

\appendix
\section{Reproducibility}
\label{sec:repro}
\sloppy
\textbf{Host.} \texttt{evo}: AMD Ryzen AI Max+ 395 (32 cores), 62.4\,GB RAM, Linux 6.17.0-35,
Go 1.26.4, glibc malloc, TeX Live 2023.\\
\textbf{Open path.} \texttt{github.com/hanzoai/dbx} \texttt{v1.16.0} +
\texttt{github.com/hanzoai/sqlite} \texttt{v0.2.1}; production driver cgo/\texttt{mattn}
(\texttt{go-sqlite3} \texttt{v1.14.47}); pure-Go control \texttt{modernc.org/sqlite}
\texttt{v1.51.0}. Pragmas exactly \texttt{sqlite.DefaultPragmas}:
\texttt{busy\_timeout=10000}, \texttt{journal\_mode=WAL},
\texttt{synchronous=NORMAL}, \texttt{foreign\_keys=ON},
\texttt{temp\_store=MEMORY}, \texttt{cache\_size=-32000}.\\
\textbf{Seed.} \texttt{base/tests/data/data.db} (344{,}064\,B) + \texttt{auxiliary.db}
(49{,}152\,B), the checked-in base-schema test database.\\
\textbf{Procedure.} Copy seed to a fresh path per tenant; open through the production path; run
the \emph{warm} workload (\texttt{count(*)} over every non-\texttt{sqlite\_} table, then
\texttt{CREATE TABLE IF NOT EXISTS} + five \texttt{INSERT}s) or the \emph{cold} workload
(\texttt{SELECT 1}); at each of $N \in \{1,10,100,1000\}$ force \texttt{runtime.GC()} +
\texttt{debug.FreeOSMemory()} and read \texttt{VmRSS} from \texttt{/proc/self/status}. Report the
$100\rightarrow1000$ incremental slope. Idle test: close all handles, re-measure; reuse test:
close all, open 1000 fresh, re-measure. Large-tenant knob: inflate the seed to 25\,MB
(60{,}000 padded rows) and repeat the \texttt{cache\_size} sweep under a full-scan warm workload.\\
\textbf{Capacity model.} Eq.~\ref{eq:model} with $U=24$\,GB, $b=0.19$\,GB, $m\in\{0.25,1.0\}$\,MB,
32\,GB nodes, HA floor 3; disk $R\times0.38$\,MB object-tiered.

\begin{thebibliography}{9}
\bibitem{hip0106} Hanzo AI. \emph{HIP-0106: Cloud --- Unified Hanzo Binary}.
\url{https://github.com/hanzoai/HIPs}.
\bibitem{unifiedbin} Z.~Kelling. \emph{The Unified Cloud Binary: Collapsing a Microservices
Control Plane into One Go Process}. Hanzo Industries white paper, 2026 (companion architecture
paper; \texttt{cloud-unified-binary}).
\bibitem{hip0116} Hanzo AI. \emph{HIP-0116: Plugin VM / In-Process Composition Model}.
\url{https://github.com/hanzoai/HIPs}.
\bibitem{hip0125} Hanzo AI. \emph{HIP-0125: ZAP-Native Transport \& gRPC Elimination} (renumbered
from HIP-0116 to avoid collision with the plugin-vm model).
\url{https://github.com/hanzoai/HIPs}.
\bibitem{hip0302} Hanzo AI. \emph{HIP-0302: Hanzo Replicate --- Encrypted SQLite Durability for
Base Services}. \url{https://github.com/hanzoai/HIPs}.
\bibitem{hip0107} Hanzo AI. \emph{HIP-0107: Replication and Durability Substrate}.
\url{https://github.com/hanzoai/HIPs}.
\bibitem{hip0032} Hanzo AI. \emph{HIP-0032: Object Storage Standard}.
\url{https://github.com/hanzoai/HIPs}.
\bibitem{hip0118} Hanzo AI. \emph{HIP-0118: SuperAdmin and Tenant Isolation Model}.
\url{https://github.com/hanzoai/HIPs}.
\bibitem{pqstore} Z.~Kelling. \emph{A Post-Quantum Object Store: Per-Tenant SQLite and a
Dropbox-Class Filesystem over One Encrypted Block Substrate}. Hanzo Industries white paper, 2026
(companion; \texttt{hanzo-pq-object-store}).
\bibitem{edgepaper} Hanzo AI. \emph{Edge LLM Inference Across Commodity Accelerators} (companion
white paper; source of the prior 154--225\,MB \texttt{cloud} working-set measurement).
\bibitem{sqlite} SQLite. \emph{PRAGMA \texttt{cache\_size}} and \emph{Write-Ahead Logging}.
\url{https://www.sqlite.org/pragma.html#pragma_cache_size}, \url{https://www.sqlite.org/wal.html}.
\bibitem{istio} Istio. \emph{Sidecar resource footprint} (Envoy proxy memory, per-proxy).
\url{https://istio.io/latest/docs/ops/deployment/performance-and-scalability/}.
\end{thebibliography}

\end{document}
